\problemname{Hemlig vandring}


I staden den exotiska och varma staden Mänhettan, bor Harry tillsammans med sin virrige kompis Tobias. 
Staden är ett $N \times N$-rutnät med koordinater $(x,y)$ där $0 \le x,y < N$. Det finns hus på alla rutor där både $x$ och $y$ är udda tal.

Tobias bor i ett av dessa hus, och hans skola ligger i ett annat. 
De ligger inte på samma rad eller kolumn, eftersom Harry har varit hemma hos Tobias och inte kunnat se skolan genom fönstret.

Varje dag går han den kortaste vägen till skolan och skickar meddelanden som ``On my way now!'' och ``Just arrived!'', 
vilket innebär att Harry vet exakt hur lång tid det tar för Tobias att röra sig från hans hus till skolan.

Trots att Harry har varit hemma hos Tobias, har Harry nu glömt var Tobias bor, men även var Tobias skola är placerad.

För att lista ut denna information, kan Harry varje dag få magiskt trolla fram väggar på valfria rutor som inte innehåller hus. 
Den dagen tar Tobias väg $k$ minuter längre än vanligt, eller $k = -1$ om han fastnar helt. Efter den dagen försvinner alla väggar som Harry hade trollat fram.

Din uppgift är att hjälpa Harry bestämma vilket hus Tobias bor i och vilket hus hans skola ligger i, med genom att spendera så få dagar som möjligt. 
Harry får spendera upp till 365 dagar på sig, men du får fler poäng ju färre dagar du använder.

\section*{Interaktion}
Det här är ett interkativt problem.

\begin{itemize}
  \item Ditt program ska först läsa in en rad med ett tal $N$ $(5 \le N \le 75, N \text{ är ett udda tal.})$, höjden och bredden av rutnätet. 
  Därefter ska ditt program läsa in $N$ rader, med $N$ karaktärer på vardera rad som beskriver rutnätet. 
  Tecknet \texttt{``\#''} innebär att den rutan är ockuperad av ett hus, medan \texttt{``.''} innebär att det inte är ockuperad av något hus.

  \item Sen ska ditt program interagera med gradern. För att 

\end{itemize}


\textbf{Se till att flusha outputen efter varje query}, annars kan du få \textit{Time Limit Exceeded}.
I C++ kan detta göras med exempelvis \texttt{cout << flush;}
eller \texttt{fflush(stdout);},
i Python med \texttt{stdout.flush()}
och i Java med \texttt{System.out.flush();}.


\section*{Poängsättning}
Din lösning kommer att testas på en mängd testfallsgrupper.
För att få poäng för en grupp så måste Harry klara alla testfall i gruppen.

\noindent
\begin{tabular}{| l | l | p{12cm} |}
  \hline
  \textbf{Grupp} & \textbf{Poäng} & \textbf{Gränser} \\ \hline
  $1$    & $14$       & $N \leq 7$ \\ \hline
  $2$    & $11$       & Tobias hus är på ruta $(1,1)$ (huset längst upp till vänster). \\ \hline
  $3$    & $\text{upp till } 75$     & Inga ytterligare begränsningar. \\ \hline
\end{tabular}

I den sista testgruppen \textbf{beror din poäng på antalet dagar som används} enligt följande formel:

\[
score = \begin{cases}
... & \text{om } 75 \le Q \le 365, \\
... & \text{om } 10 \le Q < 75, \\
... & \text{om } 6 \le Q < 10, \\
... & \text{om } Q \le 5,
\end{cases}
\]

där $Q$ är det maximala antalet tester som användes för något testfall. Poängen avrundas till närmaste heltal.